% !TeX program = xelatex
\ifdefined\BookMaster\else
\documentclass[11pt,oneside]{article}
% !TeX program = xelatex
\usepackage[b5paper,left=22mm,right=22mm,top=20mm,bottom=22mm,footskip=11mm]{geometry}
\usepackage{fontspec}
\setmainfont{TeX Gyre Termes}
\setsansfont{TeX Gyre Heros}
\usepackage{amsmath,amssymb,mathtools,amsthm}
\renewcommand{\cdots}{\mathinner{\cdotp\cdotp\cdotp}}
\usepackage{graphicx,adjustbox,array,booktabs,longtable,multirow,enumitem,changepage,needspace}
\usepackage[table]{xcolor}
\usepackage{fancyhdr}
\raggedbottom
\interfootnotelinepenalty=10000
\newcommand{\BookFootnote}[2]{\begingroup\renewcommand{\thefootnote}{\fnsymbol{footnote}}\footnote[#1]{#2}\endgroup}
\newcommand{\BookFootnoteMark}[1]{\begingroup\renewcommand{\thefootnote}{\fnsymbol{footnote}}\footnotemark[#1]\endgroup}
\newcommand{\BookFootnoteText}[2]{\begingroup\renewcommand{\thefootnote}{\fnsymbol{footnote}}\footnotetext[#1]{#2}\endgroup}
\newenvironment{BookFloat}[1]{\par\addvspace{.6\baselineskip}\noindent\def\BookFloatID{#1}\begin{minipage}{\linewidth}\phantomsection\label{bookfloat:#1:start}}{\phantomsection\label{bookfloat:\BookFloatID:end}\end{minipage}\par\addvspace{.6\baselineskip}}

\usepackage[unicode,hidelinks]{hyperref}
\pagestyle{fancy}
\fancyhf{}
\newcommand{\BookFooter}{}
\fancyfoot[L]{\parbox[b]{\dimexpr\textwidth-15mm}{\footnotesize\BookFooter}}
\fancyfoot[R]{\footnotesize\thepage}
\renewcommand{\headrulewidth}{0pt}
\renewcommand{\footrulewidth}{0pt}
\setlength{\parindent}{0pt}
\setlength{\parskip}{0.2em}
\setlength{\emergencystretch}{3em}
\setlist{nosep,leftmargin=1.7em,labelsep=.65em}
\newif\ifBookNarrative
\BookNarrativetrue
\newenvironment{BookBody}{\setlength{\parindent}{0pt}}{}
\newcommand{\BookStatement}[1]{\par\addvspace{.6\baselineskip}\noindent #1\enspace\ignorespaces}
\newcommand{\BookLabel}[1]{\par\addvspace{.6\baselineskip}\needspace{3\baselineskip}\noindent\textbf{#1}\quad\ignorespaces}
\newcommand{\BookHeading}[1]{\par\addvspace{\baselineskip}\needspace{4\baselineskip}\noindent{\large\bfseries #1}\par\medskip}
\newcommand{\BookSubheading}[1]{\par\addvspace{.7\baselineskip}\needspace{3\baselineskip}\noindent{\bfseries #1}\par\smallskip}
\newcommand{\BookTextHeading}[1]{\par\addvspace{.7\baselineskip}\needspace{3\baselineskip}\noindent #1\par\smallskip}
\newcommand{\BookSectionTitle}[2]{\par\addvspace{1.2\baselineskip}\needspace{5\baselineskip}\noindent\colorbox{gray!25}{\sffamily\bfseries\small SECTION #1}\quad{\large\bfseries #2}\par\medskip}
\newcommand{\BookExerciseTitle}[1]{\par\addvspace{\baselineskip}\needspace{4\baselineskip}\noindent{\color{gray}\rule{1.5mm}{1.5mm}}\enspace{\large\bfseries Exercises #1}\par\medskip}
\newcommand{\BookBullet}[1]{\begin{list}{\textbullet}{\setlength{\leftmargin}{1.7em}\setlength{\labelwidth}{1em}\setlength{\labelsep}{.7em}\setlength{\topsep}{0pt}\setlength{\parsep}{0pt}\setlength{\itemsep}{0pt}}\item #1\end{list}}
\newcommand{\BookExercise}[2]{\begin{list}{\textbf{#1}}{\setlength{\leftmargin}{2.2em}\setlength{\labelwidth}{1.5em}\setlength{\labelsep}{.7em}\setlength{\topsep}{.15em}\setlength{\parsep}{0pt}\setlength{\itemsep}{0pt}}\item #2\end{list}}
\newcommand{\BookEnd}[1]{\unskip\nobreak\hfill\ensuremath{#1}}
\newcommand{\BookSource}[2]{\phantomsection\label{sourcepage:#1}} % Source-page anchors, retained through reflow.
\usepackage{multicol}
\usepackage{tikz}
\newcommand{\BookPartEntry}[2]{\par\noindent\makebox[17mm][l]{\bfseries Section #1}#2\par}
\newcommand{\BookPartTitle}[3]{\par\bigskip\noindent\begin{tikzpicture}\node[anchor=north west,inner sep=0pt,outer sep=0pt] (bookpart) at (0,0) {\begin{minipage}{\dimexpr\linewidth-.5pt}\hspace*{26mm}\begin{minipage}[t]{\dimexpr\linewidth-26mm}\begin{adjustbox}{max width=\linewidth}{\fontsize{20}{24}\selectfont\bfseries #2}\end{adjustbox}\par\vspace{6mm}\hrule\vspace{6mm}{\sffamily\footnotesize\setlength{\parskip}{.3em}#3}\vspace{3mm}\end{minipage}\end{minipage}};\path[shade,top color=gray!60,bottom color=white] (0,0) rectangle ([xshift=24mm]bookpart.south west);\node[anchor=north,font=\sffamily\bfseries\small] at (12mm,-2mm) {PART};\node[anchor=north,font=\sffamily\bfseries\Huge] at (12mm,-7mm) {#1};\end{tikzpicture}\par\bigskip}
\newsavebox{\BookMathBox}
\makeatletter
\newcommand{\BookInline}[1]{\leavevmode\hskip0pt\relax\begingroup\in@{\begin}{#1}\ifin@\sbox{\BookMathBox}{\(#1\)}\ifdim\wd\BookMathBox>\dimexpr\linewidth-1em\relax\resizebox{\dimexpr\linewidth-1em\relax}{!}{\usebox{\BookMathBox}}\else\(#1\)\fi\else\in@{\left}{#1}\ifin@\sbox{\BookMathBox}{\(#1\)}\ifdim\wd\BookMathBox>\dimexpr\linewidth-1em\relax\resizebox{\dimexpr\linewidth-1em\relax}{!}{\usebox{\BookMathBox}}\else\(#1\)\fi\else\(#1\)\fi\fi\endgroup\hskip0pt\relax}
\makeatother
\DeclareMathOperator{\Aut}{Aut}
\DeclareMathOperator{\End}{End}
\DeclareMathOperator{\Hom}{Hom}
\DeclareMathOperator{\GF}{GF}
\DeclareMathOperator{\irr}{irr}
\DeclareMathOperator{\lt}{lt}
\DeclareMathOperator{\lp}{lp}
\providecommand{\textcircled}[1]{\textnormal{(#1)}}

\begin{document}
\fi
\renewcommand{\BookFooter}{Section 49 Insolvability of the Quintic}
\pdfbookmark[0]{Section 49 Insolvability of the Quintic}{49_insolvability_of_the_quintic}
\ifdefined\BookMaster\else\setcounter{page}{384}\fi
% Source PDF page 401; original printed page 384
\BookSource{401}{384}
\BookNarrativetrue

\begin{BookBody}

\BookSectionTitle{49}{Insolvability of the Quintic}

\BookTextHeading{\textbf{The Problem}  }

\noindent We are familiar with the fact that a quadratic polynomial \BookInline{f(x)= ax^{2}+ bx+ c,\allowbreak a\neq 0,\allowbreak } with real coefficients has \BookInline{(- b\pm \sqrt{b^{2}}- 4}\textit{ac}\BookInline{)/2a} as zeros in \BookInline{\mathbb{C}.} Actually, this is true for \BookInline{f(x)\in F[x],\allowbreak } where \BookInline{F} is any field of characteristic\BookInline{\neq 2} and the zeros are in \BookInline{F.} Exercise 4 asks us to show this. Thus, for example\BookInline{,\allowbreak (x^{2}+ 2x+ 3)\in \mathbb{Q}[x]} has its zeros in \BookInline{\mathbb{Q}(\sqrt{- }2).} You may wonder whether the zeros of a cubic polynomial over \BookInline{\mathbb{Q}} can also always be expressed in terms of radicals. The answer is yes, and indeed, even the zeros of a polynomial of degree 4 over \BookInline{\mathbb{Q}} can be expressed in terms of radicals. After mathematicians had tried for years to find the “radical formula” for zeros of a 5th degree polynomial, it was a triumph when Abel proved that a quintic need not be solvable by radicals. Our first job will be to describe precisely what this means. A large amount of the algebra we have developed is used in the forthcoming discussion.

\BookTextHeading{\textbf{Extensions by Radicals}  }

\BookLabel{49.1 Definition} An extension \BookInline{K} of a field \BookInline{F} is an\textbf{ extension of} \BookInline{F} \textbf{by radicals} if there are elements \BookInline{\alpha _{1},\allowbreak \cdots,\allowbreak \alpha _{r}\in K} and positive integers \BookInline{n_{1},\allowbreak \cdots,\allowbreak n_{r}} such that \BookInline{K= F(\alpha _{1},\allowbreak \cdots,\allowbreak \alpha _{r}),\allowbreak \alpha _{1}^{n_{1}}\in F} and \BookInline{\alpha _{i}^{n_{i}}\in F(\alpha _{1},\allowbreak \cdots,\allowbreak \alpha _{i- 1})} for \BookInline{1<i\leq r.} A polynomial \BookInline{f(x)\in F[x]} is\textbf{ solvable by rad}\textbf{icals over} \BookInline{F} if the splitting field \BookInline{E} of \BookInline{f(x)} over \BookInline{F} is contained in an extension of \BookInline{F} by radicals\BookInline{.}\BookEnd{\blacksquare}

A polynomial \BookInline{f(x)\in F(x)} is thus solvable by radicals over \BookInline{F} if we can obtain every zero of \BookInline{f(x)} by using a finite sequence of the operations of addition, subtraction, multiplication, division, and taking \BookInline{n_{i}}th roots, starting with elements of \BookInline{F.} Now to say that the quintic is not solvable in the classic case, that is, characteristic 0, is not to say that no quintic is solvable, as the following example shows.

\BookLabel{49.2 Example} The polynomial \BookInline{x^{5}- 1} is solvable by radicals over \BookInline{\mathbb{Q}.} The splitting field \BookInline{K} of \BookInline{x^{5}- 1} is generated over \BookInline{\mathbb{Q}} by a primitive 5th root \BookInline{\zeta } of unity. Then \BookInline{\zeta ^{5}= 1,\allowbreak } and \BookInline{K= \mathbb{Q}(\zeta ).} Similarly\BookInline{,\allowbreak x^{5}- 2} is solvable by radicals over \BookInline{\mathbb{Q},\allowbreak } for its splitting field over \BookInline{\mathbb{Q}} is generated by \BookInline{\sqrt { 2 }} and \BookInline{\zeta,\allowbreak } where \BookInline{\sqrt { 2 }} is the real zero of \BookInline{x^{5}- 2.}\BookEnd{\blacktriangle}

To say that the quintic is insolvable in the classic case means that there exists\textit{ some} polynomial of degree 5 with real coefficients that is not solvable by radicals. We shall show this.\textit{ We assume throughout this section that all fields mentioned have character}\textit{istic 0}.

The outline of the argument is as follows, and it is worthwhile to try to remember it.

\BookExercise{1.}{\textit{We shall show that} \BookInline{a} \textit{polynomial} \BookInline{f(x)\in F[x]} \textit{is solvable by radicals over} \BookInline{F(}\textit{if} \textit{and})\textit{ only if its splitting} \BookInline{\text{fi}eldE} \textit{over} \BookInline{F} \textit{has} \BookInline{a} \textit{solvable Galois group}. Recall that a solvable group is one having a composition series with\textit{ abelian} quotients. While this theorem goes both ways, we shall not prove the “if” part.}

\BookExercise{2.}{\textit{We shall show that there is} \BookInline{a} \textit{polynomial} \BookInline{f(x)\in \Phi [x]} \textit{of degree} 5\textit{ with} \BookInline{a} \textit{splitting} \BookInline{\text{fi}eldE} \textit{over} \BookInline{\mathbb{Q}} \textit{such that} \BookInline{G(E/\mathbb{Q})\simeq S_{5},\allowbreak } \textit{the symmetric group on} 5 \textit{letters}. Recall that a composition series for \BookInline{S_{5}} is \BookInline{\{\iota \}<A_{5}<S_{5}.} Since \BookInline{A_{5}} is not abelian, we will be done.}

\noindent The following lemma does most of our work for Step 1.

\BookLabel{49.3 Lemma} Let \BookInline{F} be a field of characteristic 0, and let \BookInline{a\in F.} If \BookInline{K} is the splitting field of \BookInline{x^{n}- a} over \BookInline{F,\allowbreak } then \BookInline{G(K/F)} is a solvable group.

\end{BookBody}

% Source PDF page 402; original printed page 385
\BookSource{402}{385}
\BookNarrativetrue

\begin{BookBody}

\BookTextHeading{\ensuremath{\blacksquare} \textbf{HISTORICAL NOTE}  }

\noindent The first publication of a formula for solving cubic equations in terms of radicals was in 1545 in the \BookInline{Ars} \textit{Magna} of Girolamo Cardano, although the initial discovery of the method is in part also due to Scipione del Ferro and Niccolo Tartaglia. Cardano’s student, Lodovico Ferrari, discovered a method for solving quartic equations by radicals, which also appeared in Cardano’s work.

\noindent After many mathematicians had attempted to solve quintics by similar methods, it was Joseph- Louis Lagrange who in 1770 first attempted a detailed analysis of the general principles underlying the solutions for polynomials of degree 3 and 4, and showed why these methods fail for those of higher degree. His basic insight was that in the former cases there were rational functions of the roots that took on two and three values, respectively, under all possible permutations of the roots, hence these rational functions could be written as roots of equations of degree less than that of the original. No such functions were evident in equations of higher degree.

The first mathematician to claim to have a proof of the insolvability of the quintic equation was Paolo Ruffini (1765–1822) in his algebra text of 1799. His proof was along the lines suggested by Lagrange, in that he in effect determined all of the subgroups of \BookInline{S_{5}} and showed how these subgroups acted on rational functions of the roots of the equation. Unfortunately, there were several gaps in his various published versions of the proof. It was Niels Henrik Abel who, in 1824 and 1826, published a complete proof, closing all of Ruffini’s gaps and finally settling this centuries-old question.

\BookLabel{\textit{Proof}} Suppose first that \BookInline{F} contains all the \BookInline{n}th roots of unity. By Corollary 28.7 the \BookInline{n}th roots of unity form a cyclic subgroup of \BookInline{\langle F^{\ast },\allowbreak \cdot \rangle.} Let \BookInline{\zeta } be a generator of the subgroup. (Actually, the generators are exactly the\textit{ primitive} \BookInline{n}th roots of unity.) Then the \BookInline{n}th roots of unity are

\begin{center}\begin{adjustbox}{max width=\linewidth}\BookInline{\displaystyle 1,\allowbreak \zeta ,\allowbreak \zeta ^{2},\allowbreak \cdots ,\allowbreak \zeta ^{n- 1}.}\end{adjustbox}\end{center}

\noindent If \BookInline{\beta \in F} is a zero of \BookInline{(x^{n}- a)\in F[x],\allowbreak } then all zeros of \BookInline{x^{n}- a} are

\begin{center}\begin{adjustbox}{max width=\linewidth}\BookInline{\displaystyle \beta ,\allowbreak \zeta \beta ,\allowbreak \zeta ^{2}\beta ,\allowbreak \cdots ,\allowbreak \zeta ^{n- 1}\beta .}\end{adjustbox}\end{center}

\noindent Since \BookInline{K= F(\beta ),\allowbreak } an automorphism \BookInline{\sigma } in \BookInline{G(K/F)} is determined by the value \BookInline{\sigma (\beta )} of the automorphism \BookInline{\sigma } on \BookInline{\beta.} Now if \BookInline{\sigma (\beta )= \zeta ^{i}\beta } and \BookInline{\tau (\beta )= \zeta ^{j}\beta,\allowbreak } where \BookInline{\tau \in G(K/F),\allowbreak } then

\begin{center}\begin{adjustbox}{max width=\linewidth}\BookInline{\displaystyle (\tau \sigma )(\beta )= \tau (\sigma (\beta ))= \tau (\zeta ^{i}\beta )= \zeta ^{i}\tau (\beta )= \zeta ^{i}\zeta ^{j}\beta ,\allowbreak }\end{adjustbox}\end{center}

\noindent since \BookInline{\zeta ^{i}\in F.} Similarly,

\begin{center}\begin{adjustbox}{max width=\linewidth}\BookInline{\displaystyle (\sigma \tau )(\beta )= \zeta ^{j}\zeta ^{i}\beta .}\end{adjustbox}\end{center}

\noindent Thus \BookInline{\sigma \tau = \tau \sigma,\allowbreak } and \BookInline{G(K/F)} is abelian and therefore solvable.

Now suppose that \BookInline{F} does not contain a primitive \BookInline{n}th root of unity. Let \BookInline{\zeta } be a generator of the cyclic group of \BookInline{n}th roots of unity under multiplication in \BookInline{F.} Let \BookInline{\beta } again be a zero of \BookInline{x^{n}- a.} Since \BookInline{\beta } and \BookInline{\zeta \beta } are both in the splitting field \BookInline{K} of \BookInline{x^{n}- a,\allowbreak \zeta = } \BookInline{(\zeta \beta )/\beta } is in \BookInline{K.} Let \BookInline{F^{\prime }= F(\zeta ),\allowbreak } so we have \BookInline{F<F^{\prime }\leq K.} Now \BookInline{F^{\prime }} is a normal extension of \BookInline{F,\allowbreak } since \BookInline{F^{\prime }} is the splitting field of \BookInline{x^{n}- 1.} Since \BookInline{F^{\prime }= F(\zeta ),\allowbreak } an automorphism \BookInline{\eta } in \BookInline{G(F^{\prime }/F)} is determined by \BookInline{\eta (\zeta ),\allowbreak } and we must have \BookInline{\eta (\zeta )= \zeta ^{i}} for some \BookInline{i,\allowbreak } since all zeros of \BookInline{x^{n}- 1} are powers of \BookInline{\zeta.} If \BookInline{\mu (\zeta )= \zeta ^{j}} for \BookInline{\mu \in G(F^{\prime }/F),\allowbreak } then

\begin{center}\begin{adjustbox}{max width=\linewidth}\BookInline{\displaystyle (\mu \eta )(\zeta )= \mu (\eta (\zeta ))= \mu (\zeta ^{i})= \mu (\zeta )^{i}= (\zeta ^{j})^{i}= \zeta ^{ij},\allowbreak }\end{adjustbox}\end{center}

\noindent and, similarly,

\begin{center}\begin{adjustbox}{max width=\linewidth}\BookInline{\displaystyle (\eta \mu )(\zeta )= \zeta ^{ij}.}\end{adjustbox}\end{center}

\end{BookBody}

% Source PDF page 403; original printed page 386
\BookSource{403}{386}
\BookNarrativetrue

\begin{BookBody}

\noindent Thus \BookInline{G(F^{\prime }/F)} is abelian. By Theorem 46.10,

\begin{center}\begin{adjustbox}{max width=\linewidth}\BookInline{\displaystyle \{\iota \}\leq G(K/F^{\prime })\leq G(K/F)}\end{adjustbox}\end{center}

\noindent is a normal series and hence a subnormal series of groups. The first part of the proof shows that \BookInline{G(K/F^{\prime })} is abelian, and Galois theory tells us that \BookInline{G(K/F)/G(K/F^{\prime })} is isomorphic to \BookInline{G(F^{\prime }/F),\allowbreak } which is abelian. Exercise 6 shows that if a group has a subnormal series of subgroups with abelian quotient groups, then any refinement of this series also has abelian quotient groups. Thus a composition series of \BookInline{G(K/F)} must have abelian quotient groups, so \BookInline{G(K/F)} is solvable\BookInline{.}\BookEnd{\blacklozenge}

\noindent The following theorem will complete Step 1 of our program.

\BookLabel{49.4 Theorem} Let \BookInline{F} be a field of characteristic zero, and let \BookInline{F\leq E\leq K\leq F,\allowbreak } where \BookInline{E} is a normal extension of \BookInline{F} and \BookInline{K} is an extension of \BookInline{F} by radicals. Then \BookInline{G(E/F)} is a solvable group.

\BookLabel{\textit{Proof}} We first show that \BookInline{K} is contained in a finite normal extension \BookInline{L} of \BookInline{F} by radicals and that the group \BookInline{G(L/F)} is solvable. Since \BookInline{K} is an extension of \BookInline{F} by radicals\BookInline{,\allowbreak K= } \BookInline{F(\alpha _{1},\allowbreak \cdots,\allowbreak \alpha _{r})} where \BookInline{\alpha _{i}^{n_{i}}\in F(\alpha _{1},\allowbreak \cdots,\allowbreak \alpha _{i- 1})} for \BookInline{1<i\leq r} and \BookInline{\alpha _{1}^{n_{1}}\in F.} To form \BookInline{L,\allowbreak } we first form the splitting field \BookInline{L_{1}} of \BookInline{f_{1}(x)= x^{n_{1}}- \alpha _{1}^{n_{1}}} over \BookInline{F.} Then \BookInline{L_{1}} is a normal extension of \BookInline{F,\allowbreak } and Lemma 49.3 shows that \BookInline{G(L_{1}/F)} is a solvable group. Now \BookInline{\alpha _{2}^{n_{2}}\in L_{1}} and we form the polynomial

\begin{center}\begin{adjustbox}{max width=\linewidth}\BookInline{\displaystyle f_{2}(x)= _{\sigma \in G}\prod _{(L_{1}/F)}[(x^{n_{2}}- \sigma (\alpha _{2})^{n_{2}}].}\end{adjustbox}\end{center}

\noindent Since this polynomial is invariant under action by any \BookInline{\sigma } in \BookInline{G(L_{1}/F),\allowbreak } we see that \BookInline{f_{2}(x)\in } \BookInline{F[x].} We let \BookInline{L_{2}} be the splitting field of \BookInline{f_{2}(x)} over \BookInline{L_{1}.} Then \BookInline{L_{2}} is a splitting field over \BookInline{F} also and is a normal extension of \BookInline{F} by radicals. We can form \BookInline{L_{2}} from \BookInline{L_{1}} via repeated steps as in Lemma 49.3, passing to a splitting field of \BookInline{x^{n_{2}}- \sigma (\alpha _{2})^{n_{2}}} at each step. By Lemma 49.3 and Exercise 7, we see that the Galois group over \BookInline{F} of each new extension thus formed continues to be solvable. We continue this process of forming splitting fields over \BookInline{F} in this manner: At stage \BookInline{i,\allowbreak } we form the splitting field of the polynomial

\begin{center}\begin{adjustbox}{max width=\linewidth}\BookInline{\displaystyle f_{i}(x)= _{\sigma \in G}\prod _{(L_{i- 1}/F)}[(x^{n_{i}}- \sigma (\alpha _{i})^{n_{i}}]}\end{adjustbox}\end{center}

\noindent over \BookInline{L_{i- 1}.} We finally obtain a field \BookInline{L= L_{r}} that is a normal extension of \BookInline{F} by radicals, and we see that \BookInline{G(L/F)} is a solvable group. We see from construction that \BookInline{K\leq L.}

To conclude, we need only note that by Theorem 46.8, we have \BookInline{G(E/F)\simeq G(L/F)/} \BookInline{G(L/E).} Thus \BookInline{G(E/F)} is a factor group, and hence a homomorphic image, of \BookInline{G(L/F).} Since \BookInline{G(L/F)} is solvable, Exercise 31 of Section 18 shows that \BookInline{G(E/F)} is solvable. \BookEnd{\blacklozenge}

\BookTextHeading{\textbf{The Insolvability of the Quintic}  }

\noindent It remains to find a polynomial \BookInline{f(x)\in \mathbb{Q}[x],\allowbreak } whose splitting field has Galois group \BookInline{S_{5}.} The polynomial \BookInline{f(x)= 2x^{5}- 5x^{4}+ 5} does the trick as we now show. To begin, we prove the following theorem that gives a simple condition for a subgroup of \BookInline{S_{5}} to actually be all of \BookInline{S_{5}.} We will use this theorem to show that the Galois group of the splitting field of \BookInline{f(x)} over \BookInline{\mathbb{Q}} is isomorphic with \BookInline{S_{5}.}

\BookLabel{49.5 Theorem} Let \BookInline{H} be a subgroup of \BookInline{S_{5}.} If \BookInline{H} has a transposition and a 5-cycle, then \BookInline{H= S_{5}.}

\BookLabel{\textit{Proof}} We can assume, by relabeling the points being permuted, that the 5-cycle (0, 1, 2, 3, 4) is in \BookInline{H} and the transposition \BookInline{(0,\allowbreak j),\allowbreak j\neq 0,\allowbreak } is also in \BookInline{H.} We think of \BookInline{S_{5}} as permuting the 
% Source PDF page 404; original printed page 387
\BookSource{404}{387}
elements in \BookInline{\mathbb{Z}_{5}.} By conjugating with the 5-cycle \BookInline{(0,\allowbreak 1,\allowbreak 2,\allowbreak 3,\allowbreak 4)^{r},\allowbreak } we have that for each \BookInline{0\leq r\leq 4,\allowbreak }

\end{BookBody}

\BookNarrativetrue

\begin{BookBody}

\begin{center}\begin{adjustbox}{max width=\linewidth}\BookInline{\displaystyle (0,\allowbreak 1,\allowbreak 2,\allowbreak 3,\allowbreak 4)^{r}(0,\allowbreak j)(0,\allowbreak 1,\allowbreak 2,\allowbreak 3,\allowbreak 4)^{- r}= (r,\allowbreak r+ j),\allowbreak }\end{adjustbox}\end{center}

\noindent where addition is in \BookInline{\mathbb{Z}_{5}.} Thus \BookInline{(0,\allowbreak j)} and \BookInline{(j,\allowbreak 2j)} are both in \BookInline{H;} and therefore

\begin{center}\begin{adjustbox}{max width=\linewidth}\BookInline{\displaystyle (0,\allowbreak j)(j,\allowbreak 2j)(0,\allowbreak j)= (0,\allowbreak 2j)\in H.}\end{adjustbox}\end{center}

\noindent As before, by conjugating \BookInline{(0,\allowbreak 2j)} with \BookInline{(0,\allowbreak 1,\allowbreak 2,\allowbreak 3,\allowbreak 4)^{r}} we have \BookInline{(r,\allowbreak r+ 2j)\in H} for \BookInline{0\leq r\leq 4.} Now,

\begin{center}\begin{adjustbox}{max width=\linewidth}\BookInline{\displaystyle (0,\allowbreak 2j)(2j,\allowbreak 3j)(0,\allowbreak 2j)= (0,\allowbreak 3j)\in H.}\end{adjustbox}\end{center}

\noindent Again, conjugating with powers of (0, 1, 2, 3, 4) shows that \BookInline{(r,\allowbreak r+ 3j)\in H.} Furthermore,

\begin{center}\begin{adjustbox}{max width=\linewidth}\BookInline{\displaystyle (0,\allowbreak 3j)(3j,\allowbreak 4j)(0,\allowbreak 3j)= (0,\allowbreak 4j)\in H}\end{adjustbox}\end{center}

\noindent and \BookInline{(r,\allowbreak r+ 4j)\in H.} Summarizing, we have that

\begin{center}\begin{adjustbox}{max width=\linewidth}\BookInline{\displaystyle \{(r,\allowbreak r+ sj\mid r\in \mathbb{Z}_{5}\;\text{and}\;s= \mathbb{Z}_{5}^{\ast }\}\in H.}\end{adjustbox}\end{center}

\noindent But

\begin{center}\begin{adjustbox}{max width=\linewidth}\BookInline{\displaystyle \{sj\mid s\in \mathbb{Z}_{5}^{\ast }\}= \mathbb{Z}_{5}^{\ast }}\end{adjustbox}\end{center}

\noindent since \BookInline{j\neq 0} is a unit in the field \BookInline{\mathbb{Z}_{5}.} Therefore\BookInline{,\allowbreak H} contains all the transpositions in the set

\begin{center}\begin{adjustbox}{max width=\linewidth}\BookInline{\displaystyle \{(r,\allowbreak r+ sj)\mid r\in \mathbb{Z}_{5}\;\text{and}\;s\in \mathbb{Z}_{5}^{\ast }\}= \{(r,\allowbreak r+ t)\mid r\in \mathbb{Z}_{5}\;\text{and}\;t\in \mathbb{Z}_{5}^{\ast }\}.}\end{adjustbox}\end{center}

\noindent This is the set of all transpositions in \BookInline{S_{5}.} By Theorem \BookInline{8.15,\allowbreak H= S_{5}.} \BookEnd{\blacklozenge}

\noindent We now turn our attention back to the polynomial \BookInline{f(x)= 2x^{5}- 5x^{4}+ 5.} We first observe that \BookInline{f(x)} is irreducible over \BookInline{\mathbb{Q}} by Eisenstein’s criterion, Theorem 28.16, using \BookInline{p= 5.} In order to better understand the zeros of \BookInline{f(x),\allowbreak } we compute the following.

\begin{center}\begin{adjustbox}{max width=\linewidth}\BookInline{\displaystyle \begin{aligned}f(- 1)&= - 2<0 \\ f(0)&= 5>0 \\ f(2)&= - 11<0 \\ f(3)&= 86>0\end{aligned}}\end{adjustbox}\end{center}

The intermediate value theorem from calculus says that if \BookInline{f(x)} is a continuous function and \BookInline{f(a)} and \BookInline{f(b)} have opposite signs, then \BookInline{f(x)} has a zero between \BookInline{a} and \BookInline{b.} Since the polynomial \BookInline{f(x)} is continuous, we have at least three real number zeros of \BookInline{f(x);} one between \BookInline{- 1} and 0, one between 0 and 2, and the third between 2 and 3. The derivative of \BookInline{f(x)} is

\begin{center}\begin{adjustbox}{max width=\linewidth}\BookInline{\displaystyle f^{\prime }(x)= 10x^{4}- 20x^{3}= 10x^{3}(x- 2).}\end{adjustbox}\end{center}

\noindent The only zeros of \BookInline{f^{\prime }(x)} are 0 and 2. By the mean value theorem from calculus, between any two real number zeros of \BookInline{f(x),\allowbreak } there is a zero of \BookInline{f^{\prime }(x).} Therefore, there cannot be more than three zeros of \BookInline{f(x)} that are real numbers. Let \BookInline{K} be the splitting field of \BookInline{f(x)} over \BookInline{\mathbb{Q}.} Since \BookInline{f(x)} is irreducible over \BookInline{\mathbb{Q}} and \BookInline{\mathbb{Q}} is a perfect field\BookInline{,\allowbreak K} contains five distinct zeros of \BookInline{f(x).} Thus \BookInline{f(x)} has two complex zeros, and they are complex conjugates.

We have that \BookInline{G(K/\mathbb{Q})} is isomorphic with a subgroup of the permutation group of the zeros of \BookInline{f(x).} The isomorphism maps each \BookInline{\sigma } to the permutation given by the map \BookInline{\sigma } restricted to the zeros of \BookInline{f(x).} For any zero \BookInline{\alpha \in K} of \BookInline{f(x),\allowbreak }

\begin{center}\begin{adjustbox}{max width=\linewidth}\BookInline{\displaystyle [\mathbb{Q}(\alpha ):\mathbb{Q}]= \;\text{deg}\;(\;\text{irr}\;(\alpha ,\allowbreak \mathbb{Q}))= \;\text{deg}\;(f(x))= 5.}\end{adjustbox}\end{center}

\end{BookBody}

% Source PDF page 405; original printed page 388
\BookSource{405}{388}
\BookNarrativetrue

\begin{BookBody}

\noindent Thus

\begin{center}\begin{adjustbox}{max width=\linewidth}\BookInline{\displaystyle |G(K/\mathbb{Q})|= [K:\mathbb{Q}]= [K:\mathbb{Q}(\alpha )]\cdot [\mathbb{Q}(\alpha ):\mathbb{Q}]= [K:\mathbb{Q}(\alpha )]\cdot 5.}\end{adjustbox}\end{center}

\noindent Cauchy’s Theorem 14.20 says that \BookInline{G(K/\mathbb{Q})} has a subgroup of order 5 and therefore an element of order 5. So \BookInline{G(K/\mathbb{Q})} has an element that permutes the zeros of \BookInline{f(x)} in a 5-cycle, since the only elements in \BookInline{S_{5}} that have order 5 are the 5-cycles. Also, complex conjugation is an element in \BookInline{G(K/\mathbb{Q})} that fixes each real zero of \BookInline{f(x)} and switches the two complex zeros. Thus an element of \BookInline{G(K/\mathbb{Q})} permutes the zeros of \BookInline{f(x)} by a 2-cycle. By Theorem \BookInline{49.5,\allowbreak G(K/\mathbb{Q})} is isomorphic with \BookInline{S_{5}.}

\BookLabel{49.6 Theorem} There is a polynomial in \BookInline{\mathbb{Q}[x]} that is not solvable by radicals.

\BookLabel{\textit{Proof}} As we just discovered\BookInline{,\allowbreak f(x)= 2x^{5}- 5x^{4}+ 5} is one such polynomial since its splitting field over \BookInline{\mathbb{Q}} has Galois group isomorphic with \BookInline{S_{5},\allowbreak } which is not a solvable group\BookInline{.}\BookEnd{\blacklozenge}

There are many other polynomials that are not solvable by radicals. In the proof given above, we need a few key properties for the polynomial \BookInline{f(x)} to have Galois group \BookInline{S_{5}.} First we require that \BookInline{f(x)} be an irreducible polynomial over \BookInline{\mathbb{Q}} with degree 5. We then exploit the fact that \BookInline{f(x)} has three real zeros and two complex zeros. As long as a polynomial meets these conditions, the polynomial is not solvable by radicals. The set of polynomials with these properties is infinite as shown in Exercise 8.

Exercise 9 gives a different approach to constructing polynomials in \BookInline{F[x]} whose Galois group over \BookInline{F} is \BookInline{S_{5},\allowbreak } where \BookInline{F} is a subfield of \BookInline{\mathbb{R}.} This approach also shows that there are an infinite number of polynomials that are not solvable by radicals, although the polynomials do not have rational coefficients. While Exercises 8 and 9 both produce an infinite number of polynomials that are not solvable by radicals, in Exercise 8 the number is countable while in Exercise 9 the number is uncountable.

\end{BookBody}

\BookNarrativefalse

\begin{BookBody}

\BookExerciseTitle{49}

\BookTextHeading{\textbf{Concepts}}

\BookExercise{1.}{Can the splitting field \BookInline{K} of \BookInline{x^{2}+ x+ 1} over \BookInline{\mathbb{Z}_{2}} be obtained by adjoining a square root to \BookInline{\mathbb{Z}_{2}} of an element in \BookInline{\mathbb{Z}_{2}}? Is \BookInline{K} an extension of \BookInline{\mathbb{Z}_{2}} by radicals?}

\BookExercise{2.}{Is every polynomial in \BookInline{F[x]} of the form \BookInline{ax^{8}+ bx^{6}+ cx^{4}+ dx^{2}+ e,\allowbreak } where \BookInline{a\neq 0,\allowbreak } solvable by radicals over \BookInline{F,\allowbreak } if \BookInline{F\leq \mathbb{R}}? Why or why not?}

\BookExercise{3.}{Determine whether each of the following is true or false.}

\BookExercise{a.}{Let \BookInline{F} be a field of characteristic 0. A polynomial in \BookInline{F[x]} is solvable by radicals if and only if its splitting field in \BookInline{F} is contained in an extension of \BookInline{F} by radicals.}

\BookExercise{b.}{Let \BookInline{F} be a field of characteristic 0. A polynomial in \BookInline{F[x]} is solvable by radicals if and only if its splitting field in \BookInline{F} has a solvable Galois group over \BookInline{F.}}

\BookExercise{c.}{The splitting field of \BookInline{x^{17}- 5} over \BookInline{\mathbb{Q}} has a solvable Galois group.}

\BookExercise{d.}{If \BookInline{f(x)\in \mathbb{Q}[x]} is any polynomial of degree five having three real zeros and two complex zeros, then \BookInline{f(x)} is not solvable by radicals.}

\BookExercise{e.}{The Galois group of a finite extension of a finite field is solvable.}

\BookExercise{f.}{No quintic polynomial is solvable by radicals over any field.}

\BookExercise{g.}{Every 4th degree polynomial over a field \BookInline{F\leq \mathbb{R}} is solvable by radicals.}

\BookExercise{h.}{The zeros of a cubic polynomial over a field \BookInline{F\leq \mathbb{R}} can always be attained by means of a finite sequence of operations of addition, subtraction, multiplication, division, and taking square roots starting with elements in \BookInline{F.}}

\end{BookBody}

% Source PDF page 406; original printed page 389
\BookSource{406}{389}
\BookNarrativefalse

\begin{BookBody}

\BookExercise{i.}{The zeros of a cubic polynomial over a field \BookInline{F} of characteristic 0 can never be attained by means of a finite sequence of operations of addition, subtraction, multiplication, division, and taking square roots, starting with elements in \BookInline{F.}}

\BookExercise{j.}{The theory of subnormal series of groups plays an important role in applications of Galois theory.}

\BookTextHeading{\textbf{Theory}}

\BookExercise{4.}{Let \BookInline{F} be a field, and let \BookInline{f(x)= ax^{2}+ bx+ c} be in \BookInline{F[x],\allowbreak } where \BookInline{a\neq 0.} Show that if the characteristic of \BookInline{F} is not 2, the splitting field of \BookInline{f(x)} over \BookInline{F} is \BookInline{F(\sqrt{b^{2}}- 4}\textit{ac}). [\textit{Hint:} Complete the square, just as in your high school work, to derive the “quadratic formula.”]}

\BookExercise{5.}{Show that if \BookInline{F} is a field of characteristic different from 2 and}

\begin{center}\begin{adjustbox}{max width=\linewidth}\BookInline{\displaystyle f(x)= ax^{4}+ bx^{2}+ c,\allowbreak }\end{adjustbox}\end{center}

\noindent where \BookInline{a\neq 0,\allowbreak } then \BookInline{f(x)} is solvable by radicals over \BookInline{F.}

\BookExercise{6.}{Show that for a finite group, every refinement of a subnormal series with abelian quotients also has abelian quotients, thus completing the proof of Lemma 49.3. [\textit{Hint:} Use Theorem 16.8.]}

\BookExercise{7.}{Show that for a finite group, a subnormal series with solvable quotient groups can be refined to a composition series with abelian quotients, thus completing the proof of Theorem 49.4. [\textit{Hint:} Use Theorem 16.8.]}

\BookExercise{8.}{Let \BookInline{p} be a prime number and \BookInline{f(x)= x^{5}- p^{2}x+ p\in \mathbb{Q}[x].} Prove that \BookInline{f(x)} is not solvable by radicals. [Hint: Mimic the proof that \BookInline{2x^{5}- 5x^{4}+ 5} is not solvable by radicals.]}

\BookExercise{9.}{This is an alternate method of finding a polynomial with coefficients in a field \BookInline{F,\allowbreak } a subfield of \BookInline{\mathbb{R},\allowbreak } whose Galois group over \BookInline{F} is \BookInline{S_{5}.}}

\BookExercise{a.}{Suppose that \BookInline{F\leq \mathbb{R}} is a countable field and \BookInline{x} is an indeterminate. Show that \BookInline{F[x],\allowbreak F(x),\allowbreak } and the set of real numbers that are algebraic over \BookInline{F} are all countable sets.}

\BookExercise{b.}{Show that there is a sequence of real numbers \BookInline{y_{1},\allowbreak y_{2},\allowbreak \cdots \in \mathbb{R},\allowbreak } with \BookInline{y_{1}} transcendental over \BookInline{\mathbb{Q},\allowbreak } and for each \BookInline{i>1,\allowbreak y_{i}} is transcendental over \BookInline{\mathbb{Q}(y_{1},\allowbreak y_{2},\allowbreak...,\allowbreak y_{i- 1}).}}

\BookExercise{c.}{Using the notation of part b), let \BookInline{K= \mathbb{Q}(y_{1},\allowbreak y_{2},\allowbreak y_{3},\allowbreak y_{4},\allowbreak y_{5})} and}

\begin{center}\begin{adjustbox}{max width=\linewidth}\BookInline{\displaystyle f(x)= \prod _{i= 1}^{5}(x- y_{i}).}\end{adjustbox}\end{center}

Let \BookInline{s_{1},\allowbreak s_{2},\allowbreak s_{3},\allowbreak s_{4},\allowbreak s_{5}} be the value of the elementary symmetric functions (as defined in Section 47) evaluated at \BookInline{y_{1},\allowbreak y_{2},\allowbreak y_{3},\allowbreak y_{4},\allowbreak y_{5}.} Show that \BookInline{f(x)\in F= \mathbb{Q}(s_{1},\allowbreak s_{2},\allowbreak s_{3},\allowbreak s_{4},\allowbreak s_{5}).}

\BookExercise{d.}{Show that \BookInline{K} is the splitting field of \BookInline{f(x)} over \BookInline{F} and \BookInline{G(K/F)} is isomorphic with \BookInline{S_{5}.}}

\end{BookBody}
\ifdefined\BookMaster\let\BookFinish\relax\else\def\BookFinish{\end{document}}\fi\BookFinish
